% University: Applied Science Private University

\documentclass{thesis-library}
\usepackage{array}
\renewcommand{\ttdefault}{lmtt}
\captionsetup{justification=centering}
\emergencystretch=1.5em
\newcommand{\tablefont}{\footnotesize\setstretch{1.15}\hyphenpenalty=10000\relax}

\begin{document}

\newcommand{\thesisTitle}{Library Management System (LMS)\\[0.3cm]\large Phase 1: Project Proposal \& Requirements}
\newcommand{\byLabel}{By:}
\newcommand{\studentName}{\begin{tabular}{c}Yousef Kitaneh (202311043)\\ Mohammed Abudaqqa (202310896)\\ Nasser basheer (202312323)\\ Mohammad Amin (202411102)\end{tabular}}
\newcommand{\supervisorLabel}{Instructor:}
\newcommand{\supervisorName}{Prof. Mohammad A. Shkoukani}
\newcommand{\degreeLineALabel}{Database Systems Administration}
\newcommand{\majorName}{}
\newcommand{\universityName}{Applied Science Private University}
\newcommand{\deanshipName}{Faculty of Information Technology}
\newcommand{\regionName}{Amman - Jordan}
\newcommand{\degreeyearName}{2026}
\newcommand{\degreemonthName}{November}

\makecoverpage

\pagenumbering{roman}

\newpage
\chapter*{\centering \Huge Table of Contents}
\addcontentsline{toc}{chapter}{Table of Contents}
\addtocontents{toc}{\protect\setcounter{tocdepth}{-1}}
\tableofcontents
\addtocontents{toc}{\protect\setcounter{tocdepth}{2}}

\newpage
\pagenumbering{arabic}

% ── Chapter 1: Project Proposal and Requirements Phase ──────────────────────

\printchaptertitle{Project Proposal and Requirements Phase}

\section{Chosen System}
\textbf{Library Management System (LMS)}

\section{Problem Definition}

\subsection{The Current System or Scenario}
Currently, the library operates using a fragmented, manual system. Inventory tracking relies on physical card catalogs and decentralized spreadsheets, while member registrations and daily transactions (borrowing and returning books) are recorded in paper-based ledgers. Librarians must manually check shelves or disconnected files to verify book availability and stamp due dates by hand.

\subsection{Problems that Require a Database Solution}
The existing manual setup introduces several operational bottlenecks:
\begin{itemize}
    \item \textbf{Data Redundancy and Inconsistency:} Member details and book records are duplicated across multiple sheets, leading to conflicting information.
    \item \textbf{Inefficient Search Operations:} Locating a specific book by Title, Author, or ISBN requires manual scanning, which is time-consuming and prone to human error.
    \item \textbf{Poor Transaction Tracking:} Identifying overdue books and calculating fines is a completely manual process, resulting in lost inventory and uncollected penalties.
    \item \textbf{Lack of Concurrency:} Spreadsheets do not efficiently support simultaneous multi-user access by different library staff.
\end{itemize}

\subsection{Benefits of Implementing the Database}
Transitioning to a robust Relational Database Management System (RDBMS) provides the following advantages:
\begin{itemize}
    \item \textbf{Centralized Data:} A normalized database eliminates redundancy, ensuring member and inventory data is stored once and updated globally.
    \item \textbf{Automated Tracking:} Borrowing, returning, and fine calculations are handled programmatically, ensuring accuracy and accountability.
    \item \textbf{Rapid Search and Retrieval:} Indexed tables allow instantaneous querying of the catalog, drastically reducing the time needed to serve members.
    \item \textbf{Data Integrity and Security:} The database enforces constraints to prevent invalid data entry and restricts access to sensitive member information.
\end{itemize}

\section{System Requirements}

\subsection{Functional Requirements}
The functional requirements define what the system will do:
\begin{itemize}
    \item \textbf{Add, Update, and Delete Records:} Librarians must be able to seamlessly manage book inventory (add new arrivals, update details, remove lost books) and member profiles.
    \item \textbf{Search Operations:} The system must allow users to search the library catalog by ISBN, Title, Author, Category, or Publication Year.
    \item \textbf{Manage Transactions:} The system must process book check-outs, returns, and renewals, automatically updating the available copies count.
    \item \textbf{Generate Reports:} The system must generate automated reports, including lists of currently active loans, overdue books, accumulated fines, and total inventory statistics.
\end{itemize}

\subsection{Non-Functional Requirements}
The non-functional requirements define the performance and quality standards the system must meet:
\begin{itemize}
    \item \textbf{Data Security:} The system shall implement Role-Based Access Control (RBAC). Only authorized librarians may execute modifications (INSERT, UPDATE, DELETE), while sensitive data like passwords must be hashed.
    \item \textbf{Performance:} The database must be optimized using proper indexing so that standard search queries and transaction updates execute in under two seconds.
    \item \textbf{Reliability:} The system must guarantee ACID (Atomicity, Consistency, Isolation, Durability) properties to ensure that a book check-out transaction cannot be partially completed or corrupted.
    \item \textbf{Backup and Recovery:} The database must support automated daily backups, allowing administrators to restore the system to a previous state in the event of hardware failure or accidental data loss.
\end{itemize}

\end{document}